\chapter{Metodología}

Para el desarrollo del estudio simplificado de conexión se realizaron análisis de flujo de carga y cortocircuito en el nodo 3272966 del circuito 10-502. A continuación se describen los criterios y el modelado empleados.

\section{Fuentes de Información}

\begin{itemize}
    \item Archivo \textit{Datos Circuito 10 502 - 2026.xlsx} (ESSA): parámetros de 215 líneas, 51 cargas, perfil horario de demanda y equivalente de cortocircuito.
    \item Diagrama unifilar \textit{CTO 10 502.pdf}.
    \item Fichas técnicas de inversores SOLIS 60K-LV-5G y módulos JINKO 590~W.
    \item Datos del transformador de conexión 150~kVA.
\end{itemize}

\section{Flujo de Carga}

El análisis busca determinar si, en estado estable, la incorporación del SSFV genera:

\begin{itemize}
    \item Sobrecargas en líneas y transformadores
    \item Violaciones de límites de tensión
\end{itemize}

Se aplica lo dispuesto en:
\begin{itemize}
    \item Resolución CREG 025 de 1995 (Código de Redes)
    \item Resolución CREG 174 de 2021
    \item Resolución CREG 030 de 2018
\end{itemize}

\subsection{Límites Operativos}

\begin{enumerate}
    \item \textbf{Tensiones:} aceptables entre 0,90 y 1,10~p.u. de la tensión nominal.
    \item \textbf{Cargabilidad:} hasta el 100\% de la corriente/potencia nominal de líneas y transformadores.
\end{enumerate}

\subsection{Escenarios y Franjas Horarias}

Se comparan dos escenarios (sin SSFV / con SSFV) en:
\begin{itemize}
    \item 9:00 a.m. --- demanda matutina (1,459~MVA en cabecera)
    \item 12:00 p.m. --- demanda máxima (1,683~MVA)
    \item 3:00 p.m. --- demanda vespertina (1,638~MVA)
\end{itemize}

La demanda del alimentador se ajustó a partir del registro horario del 19/03/2024 suministrado por el OR.

\section{Análisis de Cortocircuito}

Los cálculos se realizaron conforme a la norma \textbf{IEC 60909}, con tensiones de prefalla del 110\% de la tensión nominal.

\subsection{Contribuciones Consideradas}

\begin{enumerate}
    \item \textbf{Equivalente de red:} aporte de SE Conucos a 13,8~kV ($S_{k3}=654,49$~MVA; $S_{k1}=239,68$~MVA).
    \item \textbf{Inversores:} corriente de cortocircuito simétrica igual a 1,5 veces la corriente nominal del inversor (criterio habitual para esta tecnología).
\end{enumerate}

Con $S=120$~kVA y $V=220$~V:

\begin{equation}
I_{L} = \frac{120}{\sqrt{3}\times 0{,}22} = 314{,}92~\text{A}
\quad\Rightarrow\quad
I_{sc,inv}^{LV} = 1{,}5\times I_{L} = 472{,}38~\text{A}
\end{equation}

Referida a 13,2~kV:

\begin{equation}
I_{sc,inv}^{HV} = 472{,}38\times\frac{0{,}22}{13{,}2} = 7{,}87~\text{A} \approx 0{,}008~\text{kA}
\end{equation}

Este aporte es despreciable frente a los $\approx$5,12~kA trifásicos calculados en el POC.

\subsection{Parámetros Calculados}

\begin{itemize}
    \item Corriente simétrica de cortocircuito $I_k$
    \item Corriente de cresta $I_p$
\end{itemize}

\section{Modelado del Sistema}

\subsection{Red de Distribución}

Se modeló la topología completa del circuito 10-502 con:

\begin{itemize}
    \item Impedancias de secuencia positiva y cero de cada tramo
    \item Capacidades térmicas ($I_n$) por tipo de conductor
    \item Cargas LV con factor de potencia 0,9
\end{itemize}

\subsection{Transformador y Generación}

\begin{itemize}
    \item Transformador: 150~kVA, 13,2/0,22~kV, Dyn5, $Z_{cc}=4{,}9\%$
    \item Generación: 120~kW AC a FP~0,99, conectada en BT del transformador
\end{itemize}

\section{Criterios de Aceptación}

\begin{table}[H]
    \centering
    \caption{Criterios de aceptación para el análisis}
    \begin{tabular}{ll}
    \toprule
    \textbf{Parámetro} & \textbf{Criterio} \\
    \midrule
    Tensión en barras & 0,90 -- 1,10 p.u. \\
    Cargabilidad de líneas & $\leq$ 100\% \\
    Cargabilidad de transformadores & $\leq$ 100\% \\
    Aumento de cortocircuito & $<$ 1\% \\
    \bottomrule
    \end{tabular}
\end{table}
